\documentclass{ximera}

\begin{document}

    \title{Homework 3}
    
    \begin{abstract} Solve the following problems.
    \\
    Due Date: Wednesday, October 21st, 11:59 pm
    \end{abstract}
    \maketitle

\begin{exercise}
    If $P$ is an orthogonal projector, then $I-2P$ is unitary.  Prove this algebraically, and then give a geometric interpretation.    
\end{exercise}

\begin{exercise}
    Consider the matrix 
    \begin{align*}
        A = \begin{bmatrix} 1 & 2 \\ 0 & 1 \\ 1 & 0 \end{bmatrix}.
    \end{align*}
    Determine the reduced and full $QR$ factorizations $A = \hat{Q}\hat{R}$ and $A = QR$.
\end{exercise}

\begin{exercise}
    Rewrite each of the following expressions in a numerically more stable form for the indicated arguments. (Hint: trigonometric identities may come in useful.)
    \begin{enumerate}
        \item $\sqrt{x+1} -1$, $x\approx 0$
        \item $\sin(x) - \sin(y)$, $x \approx y$.
    \end{enumerate}
\end{exercise}

\begin{exercise}
    Recall that the square root operation $\sqrt{}$ represents a single floating point operation, as do addition, subtraction, multiplication, and division.  Let $A$ be an $m \times n$ matrix.  Determine the number of floating point additions, subtractions, multiplications, and divisions involved in computing the reduced QR factorization $A = \hat{Q}\hat{R}$ in Algorithm 8.1 in your textbook.
\end{exercise}

\begin{exercise}
    \begin{enumerate}
        \item Implement Algorithm 7.1 and Algorithm 8.1 from you textbook in Matlab/Python. 
        \item Consider a linearly independent list of $3$ vectors $a_1, a_2,$ and $a_3$ in $\mathbb{R}^4$.  Apply Algorithms 7.1 and 8.1 to the set of vectors which are the columns of the matrix
            $A = \begin{bmatrix}    
            1 & 1 & 1\\
            10^{-k} & 0 &0 \\
            0 & 10^{-k} & 0\\
            0 & 0 & 10^{-k}
            \end{bmatrix}
            $
            for $k = 5,6,7,8$.  
            Briefly compute the floating point algorithm error by computing the largest magnitude entry of $I - Q^TQ$, where $Q$ is a matrix where the columns are the resulting set of orthogonal vectors coming from either one of the Gram-Schmidt algorithms; record these values in a table.  Briefly explain the observed behavior in terms of numerical stability.
    \end{enumerate}
\end{exercise}


\textbf{Graduate Students:} (bonus points for undergrads)

\begin{exercise}
(Solve the following problem using Matlab or Python; as usual, submit your code to Canvas.)  

Suppose a scientific instrument records noisy observations of a
system whose underlying state lies near a two-dimensional subspace
of $\mathbb{R}^3$. We are given the following observations:
\[
X =
\begin{pmatrix}
1.3 & -0.1 & 0.4\\
0.7 & 0.766 & 0.2\\
-0.3 & 0.766 & 0.4\\
-0.9 & -0.1 & 0.2\\
-0.3 & -0.966 & 0.4\\
0.7 & -0.966 & 0.2
\end{pmatrix}.
\]

% Each row represents an observation in three-dimensional space.
% The first two coordinates approximately describe six equally
% spaced points on a circle, while the third coordinate represents
% a small perturbation.

\begin{enumerate}[label=\arabic*.]
    \item Compute the column mean vector $\bar{x}$ and form the
    centered data matrix
    \[
    X_c = X - \mathbf{1}\bar{x}^{T}.
    \]
    then compute the SVD $ X_c = U\Sigma V^T$.

    \item Construct the rank-two approximation (i.e., keep only the two leading singular values)
    \[
    X_2 = U_2\Sigma_2V_2^T.
    \]
    
    
   \textit{Fact: $X_2$ minimizes the Frobenius-norm reconstruction
    error among all matrices of rank at most two.}

    \item Calculate the residual error
    \[
    \|X_c-X_2\|_F
    \]
    and the proportion of centered data variance retained by
    the rank-two approximation.

    \item Project the centered observations onto the first two
    right singular vectors:
    \[
    Y = X_cV_2.
    \]
    Interpret the rows of $Y$ as coordinates in a
    two-dimensional principal-component space.

    \item Explain why the Euclidean distances between projected
    points equal the Euclidean distances between their rank-two
    reconstructions in the original space.
\end{enumerate}
%     \item Compute the pairwise distance matrices for the original
%     centered data and the projected data:
%     \[
%     D_{ij}=\|x_i-x_j\|_2,
%     \qquad
%     \widetilde{D}_{ij}=\|y_i-y_j\|_2.
%     \]

%     \item Establish a bound on the distortion of pairwise distances
%     introduced by projection. In particular, show that
%     \[
%     \left|
%     \|x_i-x_j\|_2-\|y_i-y_j\|_2
%     \right|
%     \leq \|r_i\|_2+\|r_j\|_2,
%     \]
%     where $r_i$ is the residual of observation $i$ after projection.
% \end{enumerate}

% \section*{Part C: Persistent Homology (25 points)}

% Using the projected data $Y$, construct the Vietoris--Rips
% filtration
% \[
% \operatorname{VR}_{\varepsilon}(Y).
% \]

% Two vertices are connected by an edge whenever their distance
% is at most $\varepsilon$. Higher-dimensional simplices are
% included whenever all their vertices are pairwise connected.

% \begin{enumerate}[label=\arabic*.]
%     \item Identify the scales at which the first edges appear
%     and the six vertices become connected into a cycle.

%     \item Determine when the one-dimensional homology class
%     $H_1$ is born and when it dies as $\varepsilon$ increases.

%     \item Calculate the corresponding Betti numbers $\beta_0$
%     and $\beta_1$ at three representative scales:
%     \begin{itemize}
%         \item Before the nearest-neighbor edges appear.
%         \item After the six-cycle forms but before it is filled
%         by triangles.
%         \item After the cycle is filled.
%     \end{itemize}

%     \item Sketch the $H_1$ persistence diagram and interpret
%     the lifetime of the loop.

%     \item Explain why a long-lived $H_1$ feature can be evidence
%     of meaningful structure, but does not by itself prove that
%     the underlying phenomenon is a true circle.
% \end{enumerate}

% \section*{Part D: The Effect of SVD Denoising on Persistence
% (20 points)}

% Suppose persistent homology is computed twice: once using the
% original centered data $X_c$ and once using the rank-two
% reconstruction $X_2$.

% \begin{enumerate}[label=\arabic*.]
%     \item Compare the two distance matrices and identify which
%     geometric features are affected by the discarded singular
%     component.

%     \item Let
%     \[
%     \delta = \max_i\|r_i\|_2.
%     \]
%     Use the stability of persistent homology to explain why
%     the persistence diagrams of the original and projected
%     datasets should be close when $\delta$ is small.

%     \item State a suitable bottleneck-distance stability bound
%     in terms of $\delta$, and explain the assumptions under
%     which it applies.

%     \item Discuss a potential failure mode: construct or describe
%     a situation in which truncating the SVD removes a
%     geometrically important feature rather than merely
%     removing noise.

%     \item Explain why selecting the truncation rank solely from
%     the percentage of variance retained may be inappropriate
%     for topological analysis.
% \end{enumerate}

% \section*{Part E: Computational Implementation (20 points)}

% Implement the analysis in Python using NumPy and a
% persistent-homology package such as GUDHI or Ripser.

% Your implementation must:
% \begin{itemize}
%     \item Center $X$ and compute its SVD.
%     \item Construct the rank-two reconstruction and projected
%     coordinates.
%     \item Compare the pairwise distance matrices.
%     \item Compute persistent homology in dimensions zero and one.
%     \item Plot the $H_1$ persistence diagrams before and after
%     projection.
%     \item Report the singular values, reconstruction error,
%     retained variance, and longest finite $H_1$ persistence
%     interval.
% \end{itemize}

\end{exercise}

\end{document}